\exercisegroup

\begin{enumerate}
\def\labelenumi{\arabic{enumi})}
\item
  设 X 是一个 Hausdorff 拓扑空间，F 是一个拓扑向量空间（在 R 或 C 上）。证明：在从 X 到 F 中的全体连续映射所组成的空间 \(\mathcal{C}(X; F)\) 上，紧收敛拓扑与向量空间结构相容。
\item
  设 \(\mathbf{E}\) 与 \(\mathbf{F}\) 是两个 Hausdorff 局部凸空间，\(\mathfrak{S}\) 是 \(\mathbf{E}\) 的有界子集的一个族。
\end{enumerate}

\begin{enumerate}
\def\labelenumi{\alph{enumi})}
\item
  证明：若 F 不只是 0，则 \(\mathcal{L}_{\mathfrak{S}}(\mathrm{E};\mathrm{F})\) 为 Hausdorff 的一个必要（且充分）条件是 S 中诸集的并在 E 中是完全的（利用 Hahn-Banach 定理）。
\item
  假设 \(\mathfrak{S}\) 是 E 的一个覆盖。证明：存在从 F 到 \(\mathcal{L}_{\mathfrak{S}}(\mathrm{E};\mathrm{F})\) 的一个闭子空间上的同构。由此推出：若 \(\mathcal{L}_{\mathfrak{S}}(\mathrm{E};\mathrm{F})\) 拟完备，则 F 必拟完备。
\item
  假设 \(\mathfrak{S}\) 是与 E 相适配有界结构（III, 第 3 页, 定义 4）。要使 \(\mathcal{L}_{\mathfrak{S}}(\mathrm{E};\mathrm{F})\) 可度量化，必须且只需 F 可度量化，且有界结构 \(\mathfrak{S}\) 具有可数基（III, 第 1 页）。要使 \(\mathfrak{S}\) -拓扑在 \(\mathcal{L}(\mathrm{E};\mathrm{F})\) 上可由单个范数定义，必须且只需 F 的拓扑可由单个范数定义，且存在一个集 \(\mathbf{M} \in \mathfrak{S}\) ，它吸收 \(\mathfrak{S}\) 中的每个集。
\end{enumerate}

\begin{enumerate}
\def\labelenumi{\arabic{enumi})}
\setcounter{enumi}{2}
\item
  设 E 是 R（相应地，C）上的一个拓扑向量空间。证明：对 R（相应地，C）的、其中没有一个是点 0 的有界子集的每个族 S，空间 \(\mathcal{L}_{\mathfrak{S}}(\mathbf{R};\mathbf{E})\) （相应地，\(\mathcal{L}_{\mathfrak{S}}(\mathbf{C};\mathbf{E})\) ）典范地同构于 E。由此推出：对每个整数 n \textgreater{} 0 以及 \(R^{n}\) （相应地，\(C^{n}\) ）的由有界子集组成的每个覆盖 S，\(\mathcal{L}_{\mathfrak{S}}(\mathbf{R}^{n};\mathbf{E})\) （相应地，\(\mathcal{L}_{\mathfrak{S}}(\mathbf{C}^{n};\mathbf{E})\) ）同构于 \(E^{n}\) 。
\item
  \begin{enumerate}
  \def\labelenumii{\alph{enumii})}
  \tightlist
  \item
    设 \(\mathrm{E}_1, \mathrm{E}_2, \mathrm{F}\) 是三个拓扑向量空间（在 \(\mathbf{R}\) 或 \(\mathbf{C}\) 上）。设 \(f\) 是从 \(\mathrm{E}_1\) 到 \(\mathrm{E}_2\) 中的连续线性映射，\(\mathfrak{S}_1\) （相应地，\(\mathfrak{S}_2\) ）是 \(\mathrm{E}_1\) （相应地，\(\mathrm{E}_2\) ）的有界子集的一个族，使得 \(f(\mathfrak{S}_1) \subset \mathfrak{S}_2\) 。证明 \(u \mapsto u \circ f\) 是从 \(\mathcal{L}_{\mathfrak{S}_2}(\mathrm{E}_2; \mathrm{F})\) 到 \(\mathcal{L}_{\mathfrak{S}_1}(\mathrm{E}_1; \mathrm{F})\) 中的一个连续线性映射。
  \end{enumerate}
\end{enumerate}

\begin{enumerate}
\def\labelenumi{\alph{enumi})}
\setcounter{enumi}{1}
\tightlist
\item
  设 \(\mathbf{E},\mathbf{F}\) 是两个拓扑向量空间，\(\mathbf{M}\) 是 \(\mathbf{E}\) 的一个向量子空间。设 \(f\) 是从 \(\mathbf{E}\) 到 \(\mathrm{E / M}\) 上的典范映射，\(\mathfrak{S}\) 是 \(\mathbf{E}\) 的有界子集的一个族。证明映射 \(u\mapsto u\circ f\) 是从 \(\mathcal{L}_{f(\mathfrak{S})}(\mathrm{E / M};\mathrm{F})\) 到 \(\mathcal{L}_{\mathfrak{S}}(\mathrm{E};\mathrm{F})\) 的如下子空间上的同构：该子空间由从 \(\mathbf{E}\) 到 \(\mathbf{F}\) 中、且在 \(\mathbf{M}\) 上为零的连续线性映射组成。
\end{enumerate}

\begin{enumerate}
\def\labelenumi{\arabic{enumi})}
\setcounter{enumi}{4}
\tightlist
\item
  设 \((\mathrm{E}_{\alpha})_{\alpha\in\mathrm{A}}\) 是一族局部凸空间，E 是一个向量空间（与诸 \(E_{\alpha}\) 有相同的标量基域），且对每个 \(\alpha\in A\) ，设 \(h_{\alpha}\) 是从 \(E_{\alpha}\) 到 E 中的一个线性映射。在 E 上赋以使诸 \(h_{\alpha}\) 连续的最细局部凸拓扑（II, 第 27 页）。对每个 \(\alpha \in \mathrm{A}\) ，设 \(\mathfrak{S}_{\alpha}\) 是 \(\mathrm{E}_{\alpha}\) 的有界子集的一个族，并设 \(\mathfrak{S}\) 是诸族 \(h_{\alpha}(\mathfrak{S}_{\alpha})\) （\(\mathrm{E}\) 的有界子集的族）的并。在这些条件下证明：对每个局部凸空间 \(\mathrm{F}\) ，\(\mathfrak{S}\) -拓扑在 \(\mathcal{L}(\mathrm{E};\mathrm{F})\) 上是使线性映射 \(u\mapsto u\circ h_{\alpha}\) （从 \(\mathcal{L}(\mathrm{E};\mathrm{F})\) 到 \(\mathcal{L}_{\mathfrak{S}_{\alpha}}(\mathrm{E}_{\alpha};\mathrm{F})\) 中）都连续的最粗拓扑。特别地，若 \(\mathrm{E}\) 是族 \((\mathrm{E}_{\alpha})_{\alpha \in \mathrm{A}}\) 的拓扑直和（II, 第 29 页, 定义 2）（每个 \(\mathrm{E}_{\alpha}\) 与 \(\mathrm{E}\) 的一个子空间等同），则乘积空间 \(\prod_{\alpha \in \mathrm{A}}\mathcal{L}_{\mathfrak{S}_{\alpha}}(\mathrm{E}_{\alpha};\mathrm{F})\) 典范地同构于空间 \(\mathcal{L}_{\mathfrak{S}}(\mathrm{E};\mathrm{F})\) ，
\end{enumerate}

其中 \(\mathfrak{S}\) 是诸 \(\mathfrak{S}_{\alpha}\) 在 \(\mathfrak{P}(\mathrm{E})\) 中的并。

\begin{enumerate}
\def\labelenumi{\arabic{enumi})}
\setcounter{enumi}{5}
\item
  设 \((\mathrm{E}_{\iota})_{\iota \in I}\) 是一族 Hausdorff 局部凸空间，其中没有一个是 0，设 E 是乘积空间 \(\prod_{\iota \in I} \mathrm{E}_{\iota}\) ，F 是一个赋范空间。证明：存在从赋以有界收敛拓扑（相应地，简单收敛拓扑、预紧收敛拓扑）的空间 \(\mathcal{L}(\mathrm{E};\mathrm{F})\) 到诸空间 \(\mathcal{L}(\mathrm{E}_{\iota};\mathrm{F})\) 的拓扑直和空间上的一个典范同构，其中每个这样的空间都赋以有界收敛拓扑（相应地，简单收敛拓扑、预紧收敛拓扑）。（注意：若 u 是从 E 到 F 中的一个连续线性映射，则存在 I 的有限子集 H，使得 \(u^{-1}(0)\) 包含诸 \(E_{\iota}\) （对所有指标 \(\iota\notin H\) ）的乘积。）
\item
  设 \(\mathrm{E}\) 、\(\mathrm{F}_1\) 、\(\mathrm{F}_2\) 是三个拓扑向量空间，\(f\) 是从 \(\mathrm{F}_1\) 到 \(\mathrm{F}_2\) 中的一个连续线性映射，\(\mathfrak{S}\) 是 \(\mathrm{E}\) 的有界子集的一个集；证明 \(u \mapsto f \circ u\) 是从 \(\mathcal{L}_{\mathfrak{S}}(\mathrm{E};\mathrm{F}_1)\) 到 \(\mathcal{L}_{\mathfrak{S}}(\mathrm{E};\mathrm{F}_2)\) 中的一个连续线性映射。
\item
  设 \(\mathbf{E}\) 是一个拓扑向量空间，带有由有界子集组成的一个集 \(\mathfrak{S}\) 。设 \((\mathbf{G}_{\iota})_{\iota \in \mathbf{I}}\) 是一族拓扑向量空间，\(\mathbf{F}\) 是一个向量空间（与 \(\mathbf{E}\) 及诸 \(\mathbf{G}_{\iota}\) 有相同的标量基域）；对每个 \(\iota \in \mathbf{I}\) ，设 \(g_{\iota}\) 是从 \(\mathbf{F}\) 到 \(\mathbf{G}_{\iota}\) 中的一个线性映射。假设在 \(\mathbf{F}\) 上赋以使诸 \(g_{\iota}\) 连续的最粗拓扑。 证明：\(\mathfrak{S}\) -拓扑在 \(\mathcal{L}(\mathbf{E};\mathbf{F})\) 上是使线性映射 \(u\mapsto g_{\iota}\circ u\) （从 \(\mathcal{L}(\mathbf{E};\mathbf{F})\) 到 \(\mathcal{L}_{\mathfrak{S}}(\mathbf{E};\mathbf{G}_{\iota})\) 中）都连续的最粗拓扑。 特别地，若 \(\mathrm{F} = \prod_{\iota \in \mathrm{I}}\mathrm{G}_{\iota}\) ，则乘积空间 \(\prod_{\iota \in \mathrm{I}}\mathcal{L}_{\mathfrak{S}}(\mathbf{E};\mathbf{G}_{\iota})\) 典范地
\end{enumerate}

等同于 \(\mathcal{L}_{\mathfrak{S}}(\mathrm{E};\mathrm{F})\)

\begin{enumerate}
\def\labelenumi{\arabic{enumi})}
\setcounter{enumi}{8}
\item
  设 E 与 F 是两个 Hausdorff 拓扑向量空间，H 是 \(\mathcal{L}(\mathrm{E};\mathrm{F})\) 的一个等度连续子集。证明：若 E 中存在一个可数的完全集，且 F 的每个有界子集都可度量化，则 H 对 E 上简单收敛拓扑可度量化。此外，若 F 的每个有界子集都满足第一可数公理，则 H 也满足第一可数公理。
\item
  设 \(\mathbf{E}\) 是一个拓扑向量空间且是 Baire 空间，\(\mathbf{F}\) 是一个拓扑向量空间。
\end{enumerate}

\begin{enumerate}
\def\labelenumi{\alph{enumi})}
\item
  证明：若一个子集 \(\mathrm{H}\) （\(\mathcal{L}(\mathrm{E};\mathrm{F})\) 的）对简单收敛拓扑有界，则 \(\mathrm{H}\) 等度连续（对 0 的每个闭邻域 \(\mathrm{V}\) （\(\mathrm{F}\) 中），考虑诸集 \(\mathbf{M}_n = \bigcap_{u\in \mathrm{H}}u^{-1}(n\mathrm{V})\) ）。
\item
  证明：若 \(\mathcal{L}(\mathrm{E},\mathrm{F})\) 的子集 H 不等度连续，则所有 \(x\in E\) （使 \(\mathrm{H}(x)\) 在 F 中不有界）组成的集是一个第一纲集的余集。由此推出：若 \((\mathrm{H}_{n})\) 是 \(\mathcal{L}(\mathrm{E};\mathrm{F})\) 的不等度连续子集的序列，则存在 \(x\in E\) ，使得诸集 \(\mathrm{H}_{n}(x)\) 中没有一个在 F 中有界（« 奇点凝聚原理 »）。
\end{enumerate}

¶ 11) 设 T 是一个可度量化的拓扑空间，E 是一个拓扑向量空间且是 Baire 空间，M 是从 E × T 到一个拓扑向量空间 F 中的映射的一个族，满足下列条件：

\(1^{\circ}\) 对一切 \(t_0 \in \mathbf{T}\) ，映射 \(x \mapsto f(x, t_0)\) （其中 \(f\) 取遍 \(\mathbf{M}\) ）全体所成的集是从 \(\mathbf{E}\) 到 \(\mathbf{F}\) 中的线性映射的一个等度连续集；

\(2^{\circ}\) 对一切 \(x_0 \in \mathbf{E}\) ，映射 \(t \mapsto f(x_0, t)\) （从 \(\mathbf{T}\) 到 \(\mathbf{F}\) 中，其中 \(f\) 取遍 \(\mathbf{M}\) ）全体所成的集是等度连续的。

证明 \(\mathbf{M}\) 等度连续。（给定 \(t_0 \in \mathbf{T}\) 以及 F 中 0 的一个闭的均衡邻域 \(\mathbf{V}\) ，对每个 \(x \in E\) ，令 \(d_{x}\) 为 T 中以 \(t_{0}\) 为中心且满足下列条件的开球半径的上确界：对这些球中任意一点 t，都有 \(f(x, t) - f(x, t_{0}) \in \mathrm{V}\) （对一切 \(f \in M\)）。证明 \(x \mapsto d_{x}\) 在每一点 \(x_{0} \in E\) 处上半连续；为此，证明：若有 \(d_{x} > \alpha > d_{x_{0}}\) （对任意接近 \(x_{0}\) 的点），那么对 F 中 0 的每个邻域 W，\(f(x_{0}, t) - f(x_{0}, t_{0})\) 都将属于 \(V + W\) （当 \(d(t, t_{0}) \leqslant \alpha\) 且 \(f \in M\) 时）。最后利用 GT, IX, § 5, No.~4, 定理 2。）

\begin{enumerate}
\def\labelenumi{\arabic{enumi})}
\setcounter{enumi}{11}
\tightlist
\item
  设 \(\mathbf{E}\) 是一个有界型局部凸空间，\(\mathfrak{S}\) 是 \(\mathbf{E}\) 的有界子集的一个族，它包含每个收敛于 0 的序列的像。
\end{enumerate}

\begin{enumerate}
\def\labelenumi{\alph{enumi})}
\item
  证明：对每个局部凸空间 \(\mathbf{F}\) ，\(\mathcal{L}_{\mathfrak{S}}(\mathrm{E};\mathrm{F})\) 的每个有界子集都等度连续。
\item
  证明：若 F 是 Hausdorff 的拟完备局部凸空间，则空间 \(\mathcal{L}_{\mathfrak{S}}(E;F)\) 是拟完备的。
\end{enumerate}

\begin{enumerate}
\def\labelenumi{\arabic{enumi})}
\setcounter{enumi}{12}
\tightlist
\item
  证明：若 E 是 Hausdorff 的半完备局部凸空间，则对每个局部凸空间 F，\(\mathcal{L}(\mathrm{E};\mathrm{F})\) 中对简单收敛拓扑有界的每个子集对每个 S-拓扑都有界。
\end{enumerate}

